\chapter{Design thinking, eller poda som flytta}

\begin{motto}
Faget eksporterte ikkje fundamentet sitt til verda.\\ Det eksporterte fråvêret av eitt, pakka som metode.
\end{motto}

Det finst eitt tilfelle der designfaget faktisk gav noko til resten av verda, og gav det med stor suksess. Dei siste tjue åra har \emph{design thinking} vandra ut av designskulane og inn i handelshøgskular, konsulenthus, departement, sjukehus og styrerom. Her, endeleg, ser det ut til at faget eksporterte ein kunnskap. Dette kapitlet viser at det motsette skjedde, og at sjølve migrasjonen — kvar poda slo rot, og kvifor — er det skarpaste enkeltbeviset i heile boka for at faget manglar ein grunn å stå på.

\section{Korleis ein prosess vart eksportvare}

Forteljinga er velkjend. Frå Stanford og designbyrået IDEO i California voks det fram, kring tusenårsskiftet, ei pakking av designpraksis som ein allmenn problemløysingsmetode. Tim Brown, leiaren i IDEO, gav han kanonisk form i \textit{Harvard Business Review} i 2008 \autocite{brown2008} og i boka \textit{Change by Design} \autocite{brown2009change}: design thinking er ein menneskesentrert prosess — empati, problemdefinisjon, idégenerering, prototyping, testing — som kven som helst i ein organisasjon kan bruke på kva som helst. David og Tom Kelley gjorde det til ei sjølvhjelpslære om «kreativt sjølvtillit» \autocite{kelley2013}. Roger Martin førte det inn i leiingsfaget som det neste konkurransefortrinnet \autocite{martin2009}. Roberto Verganti kopla det til innovasjon og meining \autocite{verganti2009}, og Jeanne Liedtka argumenterte i \textit{HBR} for at det verkar fordi det disiplinerer avgjerdsprosessar \autocite{liedtka2018}.

Suksessen var enorm, og det er verdt å spørje kvifor han var det. Svaret er ubehageleg: design thinking lét seg eksportere så lett nett fordi det var tomt der eit fag skulle vere fullt. Prosessen fortel deg å forstå brukaren, men ikkje korleis du veit at du har forstått henne, eller kva du gjer når ulike brukarar vil ha motstridande ting. Han fortel deg å generere idear, men ikkje kva som skil ein god frå ein dårleg, utover at du testar dei. Han fortel deg å lage prototypar, men ikkje kva forma faktisk følgjer, kva som verkar på henne, kvifor ho vart slik. Design thinking er ei rekkje tomme boksar med pilar mellom — og difor kan kven som helst ta han i bruk, for det krevst ingen fagkunnskap for å følgje ein tom prosedyre. Det som migrerte var ikkje ein designfagleg innsikt. Det var ein form utan innhald, og innhaldslausa var sjølve vilkåret for at han kunne reise.

\section{Kritikken innanfrå}

Dette er ikkje ein ekstern polemikars dom. Det er, i hovudsak, dommen frå designteorien sjølv, så snart han ser på fenomenet med kalde auge. Lucy Kimbell viste i to grundige artiklar at «design thinking» ikkje er eitt omgrep men fleire, dårleg definerte, som lever side om side utan å la seg sameine \autocite{kimbell2011,kimbell2012}. Johansson-Sköldberg og kollegaene hennar kartla skiljet mellom «designerly thinking» i den akademiske tradisjonen og «design thinking» i populærmanagementversjonen, og fann at den siste hadde kappa banda til den fyrste nær fullstendig \autocite{johansson2013}. Cameron Tonkinwise peikte på at populærversjonen undertrykte nett det som gjer designkompetanse til kompetanse — den situerte smaken for materiale og praksis — til fordel for ein prosess alle kunne late som dei meistra \autocite{tonkinwise2011}. Og Kees Dorst, som sjølv har arbeidd for å gje designtenking ein presis kjerne \autocite{cross2011dt}, måtte konstatere at kjernen, om han finst, ikkje er det som vart eksportert.

Frå utsida kom dommen endå hardare. Lee Vinsel kalla design thinking ein «boondoggle» — eit kostbart blendverk som lovar innovasjon og leverer post-it-lappar \autocite{vinsel2017boondoggle} — og utdjupa i \textit{The Innovation Delusion} korleis besettelsen av det nye fortrengjer det arbeidet som faktisk held samfunn oppe \autocite{vinsel2020innovation}. Natasha Iskander argumenterte i \textit{HBR} for at design thinking er grunnleggjande konservativt: det presenterer seg som radikalt, men bevarer makttilhøva ved å la dei som alt har makt definere kva som er «problemet» og kven «brukaren» er \autocite{iskander2018}. Samla teiknar denne litteraturen eit bilete faget helst ikkje ser: den mest synlege eksportartikkelen designfaget har, vart av fagfellane sine sjølve diagnostisert som innhaldslaus.

\section{Kvifor poda slo rot i økonomifaget}

Men det mest avslørande er ikkje at design thinking er tomt. Det er \emph{kvar det reiste}, og kva retninga fortel oss. Poda slo rot i økonomi- og leiingsfaget — i handelshøgskular, MBA-program, konsulentbransjen. Hadde det vore slik at design thinking var ekte designkunnskap, ville det blitt verande heime, der designkunnskap høyrer til, og blitt djupare der. I staden migrerte det, og det migrerte ikkje fordi økonomifaget hadde noko å tilby designtenkinga. Det migrerte fordi designfaget ikkje hadde nokon grunn til å halde på henne.

Tenk på kontrasten. Eit fag med eit fundament held på avkomet sitt, fordi avkomet er forankra i fagets eigen teori og veks vidare derifrå. Molekylærbiologien voks ut av biokjemien og vart verande biologi. Designfaget hadde inga teori å forankre design thinking i, ingen kumulativ tradisjon å la han modnast inni, ingen målestokk å vidareutvikle han mot. Så då handelshøgskulane tok imot — dei \emph{hadde} ein institusjonell ramme, eit språk om verdiskaping, ein marknad av betalande leiarar — fanst det ingen motkraft heime som heldt igjen. Poda dreiv dit det var jord. At jorda var økonomifaget, og ikkje designfaget sjølv, er ein dom over kvar det fanst noko å feste seg i. Designfaget var ikkje opphavet design thinking voks tilbake til; det var stasjonen han reiste frå.

Og legg merke til kva økonomifaget gjorde med han. Det stramma ikkje opp tomleiken; det utnytta han. Ein tom prosess er ideell som konsulentvare nett fordi han kan seljast til alle, ikkje kan motbevisast, og alltid kan kallast vellukka i etterkant. Design thinking vart, i si nye jord, eit instrument for å gje avgjerder som alt var tekne eit skin av brukarsentrert legitimitet. Det er ikkje ein perversjon av eit elles sunt fag. Det er kva som skjer med ein form utan innhald når han endeleg finn ein marknad: marknaden fyller han med sitt eige innhald, som her var sal.

\vignett

Migrasjonen av design thinking er faget si historie i miniatyr: ein lovnad om metode som viste seg å vere ein prosedyre utan grunn, og som difor ikkje kunne haldast heime. Men dei tomme prosessane er ikkje uskuldige. I neste del ser vi kva dei vert til når dei møter dei kraftigaste forretningsmodellane i historia — og forma vender seg mot dei ho var meint å tene.

\vidarelesing{\fullcite{brown2008}; \fullcite{kimbell2011}; \fullcite{johansson2013}; \fullcite{vinsel2020innovation}; \fullcite{iskander2018}; \fullcite{tonkinwise2011}.}
